\documentclass[10pt]{article}

% ---------- Document Identity (update these only) ----------
\newcommand{\DocStage}{}
\newcommand{\DocTitle}{Your Road to Tier One SOF}
\newcommand{\ProgramName}{182-Week Civilian-to-Selection Readiness Pipeline}
\newcommand{\ProgramSubtitle}{}

% ---------- Page ----------
\usepackage[legalpaper,left=0.5in,right=0.5in,top=0.6in,bottom=0.45in]{geometry}

% ---------- Typography ----------
\usepackage[T1]{fontenc}
\usepackage[utf8]{inputenc}
\usepackage{libertinus}
\usepackage{microtype}
\usepackage{parskip}
\usepackage{ragged2e}
\usepackage{textcomp}
\AtBeginDocument{\color{black}}
\AtBeginDocument{\pagecolor{white}}
\AtBeginDocument{\justifying}

% ---------- Landscape pages ----------
\usepackage{pdflscape}

% ---------- Color ----------
\usepackage[table]{xcolor}
\definecolor{StageBlue}{HTML}{1F4E79}
\definecolor{StageBlueDark}{HTML}{163A5A}
\definecolor{LightGray}{HTML}{F2F4F7}
\definecolor{MidGray}{HTML}{D9DEE5}
\definecolor{RuleGray}{HTML}{C7CED8}
\definecolor{HeaderInk}{HTML}{000000}
\definecolor{Ink}{HTML}{000000}

% ---------- Utility ----------
\usepackage{enumitem}
\sloppy
\setlength{\emergencystretch}{2em}

% ---------- Boxes / UI ----------
\usepackage[most]{tcolorbox}

\tcbset{
  charterTitle/.style={
    enhanced,
    colback=StageBlue!4,
    colframe=StageBlue,
    boxrule=1.25pt,
    arc=3pt,
    left=16pt,right=16pt,top=14pt,bottom=14pt,
    borderline north={3.2pt}{0pt}{StageBlue},
    borderline west={4pt}{0pt}{StageBlue},
    borderline south={1.25pt}{0pt}{StageBlue},
    drop shadow={black!25!white},
  },
  charterRule/.style={
    enhanced,
    colback=LightGray,
    colframe=RuleGray,
    boxrule=0.85pt,
    arc=3pt,
    left=12pt,right=12pt,top=10pt,bottom=10pt,
    borderline west={3pt}{0pt}{StageBlue},
  },
  charterBadge/.style={
    enhanced,
    colback=StageBlue,
    coltext=white,
    colframe=StageBlueDark,
    boxrule=0.6pt,
    arc=2pt,
    left=6pt,right=6pt,top=3pt,bottom=3pt,
  }
}
\newtcbox{\Badge}{nobeforeafter,charterBadge}

% ---------- Lists ----------
\setlist[itemize]{leftmargin=1.5em, itemsep=0.25em, topsep=0.25em}
\setlist[enumerate]{leftmargin=1.8em, itemsep=0.25em, topsep=0.25em}

% ---------- TOC ----------
\usepackage{tocloft}
\setcounter{tocdepth}{3}
\setlength{\cftbeforesecskip}{3pt}
\setlength{\cftbeforesubsecskip}{2pt}
\setlength{\cftbeforesubsubsecskip}{0pt}
\renewcommand{\cftsecfont}{\bfseries}
\renewcommand{\cftsecpagefont}{\bfseries}
\renewcommand{\cftsubsecfont}{\normalfont}
\renewcommand{\cftsubsecpagefont}{\normalfont}
\renewcommand{\cftsubsubsecfont}{\normalfont}
\renewcommand{\cftsubsubsecpagefont}{\normalfont}
\setlength{\cftsecindent}{0pt}
\setlength{\cftsubsecindent}{1.25em}
\setlength{\cftsubsubsecindent}{2.8em}
\setlength{\cftsecnumwidth}{2.1em}
\setlength{\cftsubsecnumwidth}{2.8em}
\setlength{\cftsubsubsecnumwidth}{3.6em}

% Remove the built-in TOC heading so only the custom heading is shown
\makeatletter
\renewcommand{\tableofcontents}{\@starttoc{toc}}
\makeatother

% ---------- Header/Footer: page number only ----------
\usepackage{fancyhdr}
\pagestyle{fancy}
\fancyhf{}
\renewcommand{\headrulewidth}{0pt}
\renewcommand{\footrulewidth}{0pt}
\fancyfoot[C]{\thepage}

% ---------- Section formatting ----------
\usepackage{titlesec}

\titleformat{\section}
  {\Large\bfseries\color{Ink}}
  {\thesection}{0.6em}{}
  [\vspace{0.25em}\textcolor{StageBlue}{\rule{\linewidth}{0.8pt}}]

\titleformat{\subsection}
  {\large\bfseries\color{Ink}}
  {\thesubsection}{0.6em}{}

\titleformat{\subsubsection}
  {\normalsize\bfseries\color{Ink}}
  {\thesubsubsection}{0.6em}{}

\titlespacing*{\section}{0pt}{0.95em}{0.35em}
\titlespacing*{\subsection}{0pt}{0.65em}{0.25em}
\titlespacing*{\subsubsection}{0pt}{0.55em}{0.2em}

% ---------- Table engine ----------
\usepackage{tabularray}
\UseTblrLibrary{booktabs}

% ---------- tabularray: GLOBAL table treatment ----------
\SetTblrInner{
  cells = {fg=black, bg=white, valign=t},
}

% ---------- Header helpers ----------
\newcommand{\Hdr}[1]{\shortstack[c]{\strut #1\strut}}

% ---------- Lines ----------
\newcommand{\TitleLine}{\textcolor{StageBlue}{\rule{0.98\linewidth}{0.95pt}}}
\newcommand{\SubTitleLine}{\textcolor{RuleGray}{\rule{0.98\linewidth}{0.65pt}}}

% ---------- Continued on next page ----------
\newcommand{\ContinuedOnNextPageInline}{%
  \par
  \vfill
  \noindent\textcolor{RuleGray}{\rule{\linewidth}{1pt}}\vspace{-2pt}
  \noindent\hfill\textit{Continued on next page}\par\vspace{-10pt}
  \noindent\textcolor{RuleGray}{\rule{\linewidth}{1pt}}
  \clearpage
}

% ---------- PDF hyperlinks ----------
\usepackage[hidelinks]{hyperref}
\hypersetup{
  colorlinks=false,
  pdfborder={0 0 0},
  linktoc=all,
  pdftitle={\DocTitle},
  pdfauthor={\ProgramName}
}

% =========================================================
\begin{document}

\section*{Stage 2.D — Performance Thresholds and Alignment for Entry, Intermediate, and Selection-Ready}
\phantomsection
\addcontentsline{toc}{section}{Stage 2.D — Performance Thresholds and Alignment for Entry, Intermediate, and Selection-Ready}



\subsection*{2.D.1 Tier Definitions}
\phantomsection
\addcontentsline{toc}{subsection}{2.D.1 Tier Definitions}

\textbf{ENTRY} tier (governance definition):

\begin{itemize}
\item \textbf{ENTRY} \textbf{MUST} be treated as minimum safe participation and baseline admissibility under \textbf{VALID} conditions.
\item \textbf{ENTRY} indicates the candidate can complete baseline test constructs without violating safety posture and without producing invalid evidence.
\item \textbf{ENTRY} is not “ready” and \textbf{MUST} NOT be used to justify early specificity escalation, impact escalation, or aggressive intensity increases.
\end{itemize}

\textbf{INTERMEDIATE} tier (governance definition):

\begin{itemize}
\item \textbf{INTERMEDIATE} \textbf{MUST} represent stabilized tolerance and progressing capacity under \textbf{VALID} conditions.
\item \textbf{INTERMEDIATE} indicates performance improvement combined with sufficient durability and recovery stability to support higher specificity and controlled intensity increases without predictably triggering breakdown.
\item \textbf{INTERMEDIATE} classification \textbf{MUST} remain constrained by safety governance (stop rules, pain-altered mechanics rules, water governance) and \textbf{MUST} NOT be treated as permission to “push through” risk flags.
\end{itemize}

\textbf{SELECTION-READY} tier (governance definition):

\begin{itemize}
\item \textbf{SELECTION-READY} \textbf{MUST} align to the Stage 0.3 end-state envelope and \textbf{MUST} be interpreted as a converged system state across domains, not a single “hit it once” outcome.
\item \textbf{SELECTION-READY} tier attainment \textbf{MUST} be supported by repeatable, \textbf{VALID} evidence and stable durability/recovery posture.
\end{itemize}

Replication rule (locked; applies to any tier classification used in gates):

\begin{itemize}
\item A tier classification used for progression decisions \textbf{MUST} be supported by stability: at least 2 of the last 3 \textbf{VALID} test events within the decision window for that metric.
\item If replication is not met, the result MAY be recorded, but tier classification defaults to \textbf{HOLD} for progression decisions for that metric until replication exists.
\end{itemize}

Evidence rule (locked):

\begin{itemize}
\item Thresholds apply only to \textbf{VALID} tests.
\item \textbf{INVALID} tests do not qualify for tier attainment and provide no gate credit.
\end{itemize}

\subsubsection*{2.D.1.1 Tier Classification Is Conditional on Admissibility and Comparability}
\phantomsection
\begin{itemize}
\item Tier classification is not assigned from results that cannot be verified under Stage 2 validity conditions (\textbf{ENV}/\textbf{EQUIP}/\textbf{WARMUP}/\textbf{SAFETY}/\textbf{MEASURE}) and Stage 1.F admissibility posture.
\item Tier classification is not assigned across a comparability break unless the comparability break is explicitly logged and the evidence is repeated under comparable conditions.
\item Tier classification is not assigned when disqualifier controls are active (for example, gait-altering pain flags, active stop-rule events, or underwater governance failures).
\end{itemize}

\subsubsection*{2.D.1.2 Metric Posture States for Gate Use}
\phantomsection

For each Metric \textbf{ID} in the threshold table, the gate interpretation posture \textbf{MUST} be one of the following:

\begin{itemize}
\item \textbf{UNCLASSIFIED} — no \textbf{VALID} evidence exists in the decision window for the metric (missing or invalid tests; missing trend evidence).
\item \textbf{CLASSIFIED} (\textbf{ENTRY}) — meets \textbf{ENTRY} threshold under replication rule and admissibility rules.
\item \textbf{CLASSIFIED} (\textbf{INTERMEDIATE}) — meets \textbf{INTERMEDIATE} threshold under replication rule and admissibility rules.
\item \textbf{CLASSIFIED} (\textbf{SELECTION-READY}) — meets \textbf{SELECTION-READY} threshold under replication rule and admissibility rules.
\item \textbf{HOLD} (Replication) — a result exists but replication rule (2 of last 3 \textbf{VALID}) is not met.
\item \textbf{HOLD} (Disqualifier) — replication exists but disqualifier controls constrain tier use for progression decisions.
\item NOT \textbf{APPLICABLE} — metric is not applicable by design for the current phase or gate context (used only when the metric is explicitly out-of-scope for the gate).
\end{itemize}
\clearpage
\subsection*{2.D.2 Threshold Schema and Rules}
\phantomsection
\addcontentsline{toc}{subsection}{2.D.2 Threshold Schema and Rules}

Threshold applicability (locked):

\begin{itemize}
\item Thresholds apply only to \textbf{VALID} tests. Any threshold claim based on \textbf{INVALID} evidence is inadmissible.
\item If a test is missing or invalid, no tier classification is assigned for that metric; the metric posture is \textbf{UNCLASSIFIED} (no valid data) for gate use.
\end{itemize}

Directionality definitions (locked):

\begin{itemize}
\item Lower is better: time metrics (\textbf{RUN}, \textbf{RUCK}, \textbf{SWIM}, \textbf{AER} time trials, underwater time where governed).
\item Higher is better: reps (\textbf{CAL}) and holds (\textbf{CAL} hold time in seconds).
\item Pass band: \textbf{BODYCOMP} envelope targets (bodyweight and body fat range).
\item Binary: \textbf{COMPLIANCE} and safety flags (Y/N; or enumerated tags where defined).
\end{itemize}

Service-test anchoring rule (locked posture):

\begin{itemize}
\item \textbf{ENTRY} and \textbf{INTERMEDIATE} tiers \textbf{SHOULD} be anchored to published minimums and competitive equivalents where the metric exists in common feeder tests (\textbf{CAL} and common timed runs).
\item Where no direct feeder-test analog exists, \textbf{ENTRY} and \textbf{INTERMEDIATE} thresholds \textbf{MUST} be derived conservatively from \textbf{SELECTION-READY} anchors and safety posture, and \textbf{MUST} NOT imply early readiness.
\item Any derived threshold \textbf{MUST} include a brief derivation basis statement in Notes in the master table.
\end{itemize}

Replication rule (explicit and numeric; locked):

\begin{itemize}
\item Tier classification used for gate decisions requires at least 2 of last 3 \textbf{VALID} test events within the decision window.
\end{itemize}

Use posture and crosswalk posture (locked):

\begin{itemize}
\item Phase Rows import thresholds by Metric \textbf{ID} and reference Protocol Cards where applicable.
\item Stage 3 schedules test cadence and deload placement; this deliverable defines tier numbers only and does not place tests on a calendar.
\end{itemize}

Inputs and Assumptions (required)

Inputs (Binding):

\begin{itemize}
\item Stage 0.3 Selection-Ready anchors (bodyweight/body fat; calisthenics; running; swimming; rucking; underwater time range framing under governance).
\item Stage 2.A canonical domain tags.
\item Stage 2.B metric IDs (\textbf{MET-2B}-\#\#\#).
\item Stage 2.C protocol cards (\textbf{C1}–\textbf{C18}), where applicable.
\item Replication rule: 2 of last 3 \textbf{VALID} tests required for tier classification used in gates.
\item Validity posture: thresholds apply only to \textbf{VALID} tests; invalid tests yield no gate credit.
\end{itemize}

Assumption Ledger (minimum; required due to \textbf{ID}/protocol coverage gaps):

\begin{itemize}
\item \textbf{MET-2B-TBD-001} and \textbf{MET-2B-TBD-002} are used to express required session validity distribution posture as measurable compliance thresholds (weekly \textbf{MODIFIED} count and weekly \textbf{INVALID} count). These placeholders \textbf{MUST} be reconciled into Stage 2.B before Stage 2 is considered complete.
\item \textbf{MET-2B-018} through \textbf{MET-2B-021} are covered by \textbf{C12}. No separate \textbf{TBD} protocol card is required; gate use depends on compliant \textbf{C12} execution.
\item \textbf{MET-2B-032} through \textbf{MET-2B-034} are covered by \textbf{C18}. Numeric thresholds remain supporting only until \textbf{C18} is finalized and its placeholder sections are fully resolved.
\item Strength proxy thresholds (\textbf{MET-2B-027} through \textbf{MET-2B-031}) are not anchored in Stage 0.3. Numeric thresholds are provided as conservative governance thresholds subject to versioned refinement once the fixed construct and implement posture are stabilized in protocol.
\end{itemize}

Reconciliation Notes (conflict resolved; adopted posture applied here):

\begin{itemize}
\item Legacy materials use non-stage-scoped metric IDs (e.g., \textbf{MET}-\#\#\#) and differing protocol card numbering for some constructs. This deliverable adopts the locked conventions for this package: metric IDs are \textbf{MET-2B}-\#\#\# and protocol cards are \textbf{C1}–\textbf{C18}. Any legacy \textbf{ID} references are treated as non-authoritative identifiers and must be mapped in the crosswalk layer; downstream phases \textbf{MUST} reference \textbf{MET-2B}-\#\#\# only.
\end{itemize}

\subsubsection*{2.D.2.1 Decision Window Definition}
\phantomsection

\begin{itemize}
\item The “decision window” is the gate’s evidence window as defined by Stage 1.F gate doctrine and Stage 3 calendar posture.
\item Thresholds and replication are evaluated within the decision window only.
\item A test event outside the decision window may remain in the historical record but cannot be used to satisfy “2 of last 3 \textbf{VALID} within the decision window” unless the decision window explicitly includes it.
\end{itemize}

\subsubsection*{2.D.2.2 “Last 3 \textbf{VALID} Tests” Definition }
\phantomsection

\begin{itemize}
\item “Last 3 \textbf{VALID} test events” refers to the three most recent test events for the metric that are tagged \textbf{VALID} and fall within the decision window.
\item If fewer than three \textbf{VALID} test events exist within the decision window, use all available \textbf{VALID} events for replication evaluation and apply the replication rule conservatively:
  \begin{itemize}
  \item 0 \textbf{VALID} → \textbf{UNCLASSIFIED}.
  \item 1 \textbf{VALID} → \textbf{HOLD} (Replication).
  \item 2 \textbf{VALID} → tier may be classified if both meet threshold; if mixed, classify to the lower tier or \textbf{HOLD} depending on gate posture.
  \item 3+ \textbf{VALID} → apply “2 of last 3” rule.
  \end{itemize}
\end{itemize}

\subsubsection*{2.D.2.3 Comparability Preconditions for Applying Thresholds}
\phantomsection

Threshold claims are admissible only when the test events are comparable under recorded conditions:

\begin{itemize}
\item ROCK is not used; the canonical tag is \textbf{RUCK}.
\item \textbf{RUN} metrics: measured route or treadmill allowed only when the protocol card validity fields are satisfied; treadmill grade and settings must be stable across the events used for replication.
\item \textbf{RUCK} metrics: dry load only; load verification and configuration must be stable across replication events; route/machine identity must be stable.
\item \textbf{SWIM} metrics: pool unit (yards vs meters) must match; allowable strokes and no fins posture must be satisfied; pool length must be verified consistently.
\item \textbf{BODYCOMP} metrics: method consistency within the window is required for trend interpretability; method drift requires comparability break handling.
\item \textbf{AER} metrics: construct definition (distance or duration) and machine settings must be fixed in protocol before tiering is gate-admissible.
\end{itemize}

If comparability is materially broken across replication events, tier classification defaults to \textbf{HOLD} until repeated under comparable conditions.
\newpage
\begin{landscape}
\subsection*{2.D.3 Master Threshold Table}
\phantomsection
\addcontentsline{toc}{subsection}{2.D.3 Master Threshold Table}

\begingroup
\scriptsize
\setlength{\tabcolsep}{2pt}
\renewcommand{\arraystretch}{1.12}

\begin{longtblr}{
  width=\linewidth,
  rowhead=1,
  colspec = {
    X[l,0.85]
    X[l,1.50]
    X[l,0.75]
    X[l,0.75]
    X[l,1.15]
    X[l,0.95]
    X[l,1.10]
    X[l,1.35]
    X[l,2.80]
    X[l,1.05]
  },
  hlines,
  vlines,
  cells = {hsep=6pt, vsep=6pt, halign=l, valign=t},
  cell{1-Z}{1-10} = {fg=black},
  row{odd} = {bg=white},
  row{even} = {bg=LightGray},
  row{1} = {bg=StageBlue!15, fg=black, font=\bfseries, halign=l, valign=t},
}
\Hdr{Metric \textbf{ID}} &
\Hdr{Metric Name} &
\Hdr{Domain Tag} &
\Hdr{Protocol Card} &
\Hdr{Directionality} &
\Hdr{Entry Threshold} &
\Hdr{Intermediate Threshold} &
\Hdr{Selection-Ready Threshold} &
\Hdr{Notes} &
\Hdr{Overbuild Preferred} \\
\textbf{MET-2B-001} & Bodyweight — Weekly Average & \textbf{BODYCOMP} & \textbf{C1} & Pass band (lb) & ≤260 lb (weekly average) & ≤220 lb (weekly average) & 175–185 lb (center-point 180 lb) & \textbf{ENTRY}/\textbf{INTERMEDIATE} are conservative staging anchors to prevent premature readiness claims; Selection-Ready is expressed as a pass band around the ~180 lb center-point; weekly average required (not single weigh-in) & 175–180 lb (weekly average) \\
\textbf{MET-2B-002} & Body Fat Estimate Percent — Method/Device Recorded & \textbf{BODYCOMP} & \textbf{C2} & Pass band (\%) & ≤40\% (method recorded) & ≤25\% (method consistent within window) & 12–17.5\% & Method-agnostic at governance level; trend validity requires consistent method across measurement windows; Selection-Ready range matches end-state envelope & 12–15\% \\
\textbf{MET-2B-010} & Push-Ups — 2-Minute Timed, Strict & \textbf{CAL} & \textbf{C4} & Higher is better (reps) & ≥25 reps & ≥60 reps & ≥100 reps & \textbf{ENTRY}/\textbf{INTERMEDIATE} are conservative placeholders to stage progression safely; \textbf{VALID} gate use requires compliant \textbf{CAL} artifact evidence and replication posture & ≥110 reps \\
\textbf{MET-2B-011} & Sit-Ups — 2-Minute Timed, Standard & \textbf{CAL} & \textbf{C5} & Higher is better (reps) & ≥25 reps & ≥60 reps & ≥100 reps & APPLY ONLY TO \textbf{VALID} tests; do not tier from invalid counting conditions; replication required for gate classification & ≥110 reps \\
\textbf{MET-2B-012} & Pull-Ups — Strict Dead-Hang & \textbf{CAL} & \textbf{C6} & Higher is better (reps) & ≥3 reps & ≥10 reps & ≥20 reps & Pronated strict dead-hang standard governs validity; missing required artifact evidence yields \textbf{INVALID} for gate use & ≥22 reps \\
\textbf{MET-2B-013} & Plank — Forearm & \textbf{CAL} & \textbf{C7} & Higher is better (seconds) & ≥120 s (2:00) & ≥180 s (3:00) & ≥225 s (3:45) & Termination on form failure; validity requires standardized setup and auditable timing; replication required & ≥240 s (4:00) \\
\textbf{MET-2B-014} & 1.5-Mile Run Time Trial & \textbf{RUN} & \textbf{C8} & Lower is better (mm:ss) & ≤15:00 & ≤12:00 & ≤9:00 & Impact-managed; measured route OR treadmill permitted when declared and logged; \textbf{ENTRY}/\textbf{INTERMEDIATE} are staging markers and do not authorize unsafe impact escalation & ≤8:30 \\
\textbf{MET-2B-015} & 2-Mile Run Time Trial & \textbf{RUN} & \textbf{C9} & Lower is better (mm:ss) & ≤20:00 & ≤16:00 & ≤12:30 & Validity requires measured distance or treadmill settings logged; do not tier from \textbf{INVALID} conditions; replication required & ≤12:00 \\
\textbf{MET-2B-016} & 3-Mile Run Time Trial & \textbf{RUN} & \textbf{C10} & Lower is better (mm:ss) & ≤33:00 & ≤27:00 & ≤20:00 & Entry/Intermediate are conservative staging markers; phases govern impact exposure ramps and fatigue control & ≤19:00 \\
\textbf{MET-2B-017} & 5-Mile Run Time Trial & \textbf{RUN} & \textbf{C11} & Lower is better (mm:ss) & ≤60:00 & ≤45:00 & ≤35:00 & Selection-Ready threshold is 35:00 or faster; longer duration increases heat/hydration confounders; conditions logging required & ≤30:00 \\
\textbf{MET-2B-022} & Swim 500 yd Time — No Fins & \textbf{SWIM} & \textbf{C13} & Lower is better (mm:ss) & ≤20:00 & ≤14:00 & ≤10:30 & No fins; allowable strokes per protocol; pool length must be verified; lane interruptions can invalidate; replication required & ≤9:30 \\
\textbf{MET-2B-023} & Swim 500 m Time — No Fins & \textbf{SWIM} & \textbf{C13} & Lower is better (mm:ss) & ≤20:00 & ≤14:00 & ≤10:00 & No fins; allowable strokes per protocol; meters unit is locked; no conversions without mode declaration; replication required & ≤9:00 \\
\textbf{MET-2B-025} & Underwater Exposure Admissibility Flag & \textbf{SWIM} & \textbf{C14} & Binary (Y/N) & N/A (domain remains prohibited unless admissibility passes) & Y required before any underwater testing & Y required before any underwater performance metric & Governance gate, not performance target; if “N,” underwater remains prohibited; any governance breach invalidates readiness claims and forces conservative posture & Y (maintained across decision window) \\
\textbf{MET-2B-026} & Underwater 25 m Time — One Breath & \textbf{SWIM} & \textbf{C14} & Lower is better (seconds) & Not required at \textbf{ENTRY}; governed domain locked out until admissibility passes & Not required at \textbf{INTERMEDIATE}; governance prerequisites dominate & 30–45 s (controlled) & Performance threshold applies only after admissibility passes (\textbf{MET-2B-025} = Y) and under certified supervision + safety spotter; thresholds \textbf{MUST} NOT be interpreted as encouragement for unsupervised practice & <30 s (controlled; not required) \\
\textbf{MET-2B-018} & 3-Mile Ruck Time Trial — 25 lb (dry) & \textbf{RUCK} & \textbf{C12} & Lower is better (mm:ss) & ≤75:00 & ≤60:00 & ≤45:00 & Dry load only; water carried logged separately; treadmill allowed only with Environment Unit \textbf{ID} and settings logged; \textbf{ENTRY}/\textbf{INTERMEDIATE} use looser times to avoid early heavy-load speed demands & ≤42:00 \\
\textbf{MET-2B-019} & 6-Mile Ruck Time Trial — 35 lb (dry) & \textbf{RUCK} & \textbf{C12} & Lower is better (mm:ss) & ≤150:00 & ≤120:00 & ≤90:00 & Dry load only; water carried logged separately; validity requires load verification and stable route/tools; looser Entry/Intermediate times prevent premature speed/load escalation & ≤85:00 \\
\textbf{MET-2B-020} & 9-Mile Ruck Time Trial — 45 lb (dry) & \textbf{RUCK} & \textbf{C12} & Lower is better (mm:ss) & ≤240:00 & ≤195:00 & ≤135:00 & \textbf{C12} governs this metric; dry load only; water logged separately; times are staged to avoid early heavy-load speed & ≤125:00 \\
\textbf{MET-2B-021} & 12-Mile Ruck Time Trial — 55 lb (dry) & \textbf{RUCK} & \textbf{C12} & Lower is better (mm:ss) & ≤300:00 & ≤240:00 & ≤180:00 & \textbf{C12} governs this metric; dry load only; water logged separately; selection-ready aligns to 15:00/mi pace standard & ≤165:00 \\
\textbf{MET-2B-046} & Weekly Admissible Session Compliance Percent & \textbf{RECOVERY} & N/A & Higher is better (\%) & ≥67\% (≥4 of 6 admissible, average across decision window) & ≥83\% (≥5 of 6 admissible, average across decision window) & ≥83\% (≥5 of 6 admissible, average across decision window) & Admissible = \textbf{VALID} + admissible \textbf{MODIFIED} within cap posture; \textbf{INVALID} is non-admissible; thresholds apply to the decision window used by the gate; replication posture does not apply (weekly compliance is a governance statistic) & ≥92\% (≥11 of 12 across a 2-week window) \\
\textbf{MET-2B-TBD-001} & Weekly \textbf{MODIFIED} Session Count (non-medical) & \textbf{RECOVERY} & N/A & Lower is better (count) & ≤2 per week & ≤1 per week & ≤1 per week & Expresses session validity distribution posture; outside medical issues, \textbf{MODIFIED} sessions must be capped; repeated \textbf{MODIFIED} weeks force conservative gate posture even if test performance improves & 0 per week \\
\textbf{MET-2B-TBD-002} & Weekly \textbf{INVALID} Session Count & \textbf{RECOVERY} & N/A & Lower is better (count) & ≤1 per week & 0 per week & 0 per week & Expresses session validity distribution posture; \textbf{INVALID} sessions are non-admissible and break gate eligibility; any sustained \textbf{INVALID} pattern forces \textbf{HOLD}/\textbf{REGRESS} posture & 0 per week (maintained) \\
\textbf{MET-2B-045} & Cardio Minimum Standard Met (Decision Window) & \textbf{RECOVERY} & N/A & Higher is better (\%) & ≥83\% of sessions in window & 100\% of sessions in window & 100\% of sessions in window & This is a compliance metric; “met” means cardio minimum standard satisfied per session; failure reduces admissible session count and degrades gate eligibility & 100\% with stable modality logging \\
\textbf{MET-2B-036} & Gait-Altering Pain Flag & \textbf{DURABILITY} & N/A & Binary (Y/N) & No during any \textbf{VALID} test event; if Yes in window, tier classification defaults \textbf{HOLD} for progression decisions & No & No & Binary pass condition for tier classification used in progression: must be No at \textbf{INTERMEDIATE} and \textbf{SELECTION-READY}; any Yes triggers conservative posture and invalidates escalation claims & No (maintained) \\
\textbf{MET-2B-027} & Hinge Strength Proxy — \textbf{3-RM} or \textbf{5-RM} Load (submax) & \textbf{STR} & \textbf{C17} & Higher is better (lb) & ≥135 lb & ≥225 lb & ≥315 lb & Governance thresholds (not Stage 0.3 anchored); evaluated only under submax proxy posture (no \textbf{1-RM}); implement type and load method must be held stable within verification windows & ≥365 lb \\
\textbf{MET-2B-028} & Squat Strength Proxy — \textbf{3-RM} or \textbf{5-RM} Load (submax) & \textbf{STR} & \textbf{C17} & Higher is better (lb) & ≥115 lb & ≥205 lb & ≥275 lb & Governance thresholds (not Stage 0.3 anchored); submax proxy only; validity requires implement identity and standardized depth/controls per protocol & ≥315 lb \\
\textbf{MET-2B-029} & Upper Push Strength Proxy — \textbf{3-RM} or \textbf{5-RM} Load (submax) & \textbf{STR} & \textbf{C17} & Higher is better (lb) & ≥75 lb & ≥135 lb & ≥185 lb & Governance thresholds (not Stage 0.3 anchored); submax proxy only; implement identity and rep scheme must be logged and held stable within verification windows & ≥205 lb \\
\textbf{MET-2B-030} & Upper Pull Strength Proxy — \textbf{3-RM} or \textbf{5-RM} Load (submax) & \textbf{STR} & \textbf{C17} & Higher is better (lb) & ≥80 lb & ≥150 lb & ≥200 lb & Governance thresholds (not Stage 0.3 anchored); submax proxy only; validity requires stable implement posture and logged method & ≥225 lb \\
\textbf{MET-2B-031} & Carry Capacity — Load Outcome (submax proxy) & \textbf{STR} & \textbf{C17} & Higher is better (lb) & ≥150 lb & ≥250 lb & ≥350 lb & Governance thresholds (not Stage 0.3 anchored); carry distance/time construct must be fixed in protocol; until fixed, treat tiering as provisional and interpret conservatively & ≥400 lb \\
\textbf{MET-2B-032} & Rower Time Trial & \textbf{AER} & \textbf{C18} & Lower is better (mm:ss) & ≤12:00 & ≤10:00 & ≤8:30 & Numeric thresholds remain supporting until \textbf{C18} is finalized and non-placeholder; \textbf{AER} is not interchangeable evidence for \textbf{RUN} readiness & ≤8:00 \\
\textbf{MET-2B-033} & Bike Time Trial & \textbf{AER} & \textbf{C18} & Lower is better (mm:ss) & ≤30:00 & ≤25:00 & ≤22:00 & Numeric thresholds remain supporting until \textbf{C18} is finalized and non-placeholder; do not use \textbf{AER} tier attainment to justify impact escalation & ≤20:00 \\
\textbf{MET-2B-034} & Elliptical Time Trial & \textbf{AER} & \textbf{C18} & Lower is better (mm:ss) & ≤20:00 & ≤16:00 & ≤14:00 & Numeric thresholds remain supporting until \textbf{C18} is finalized and non-placeholder; \textbf{AER} remains non-impact conditioning domain & ≤13:00 \\
\end{longtblr}

\end{landscape}
\endgroup
\subsection*{2.D.4 Alignment Notes}
\phantomsection
\addcontentsline{toc}{subsection}{2.D.4 Alignment Notes}

How thresholds are consumed downstream:

\begin{itemize}
\item Phase Rows (Stages 4–5) reference Metric IDs (\textbf{MET-2B}-\#\#\#) and tier thresholds (\textbf{ENTRY} / \textbf{INTERMEDIATE} / \textbf{SELECTION-READY}) as gate context. Phase Rows do not rewrite thresholds; they import them.
\item Protocol Cards (Stage 2.C) define how tests are executed and what makes them \textbf{VALID} or \textbf{INVALID}. Thresholds have no operational meaning without \textbf{VALID} tests.
\item Stage 3 schedules testing and deload placement at a high level and preserves validity by controlling fatigue confounders and reducing novelty near verification windows. This deliverable does not place tests on a calendar.
\end{itemize}

How tiers prevent premature escalation:

\begin{itemize}
\item \textbf{RUN} thresholds do not authorize unsafe impact progression. Meeting an \textbf{ENTRY}/\textbf{INTERMEDIATE} run time does not override pain-altered mechanics rules, durability flags, or environmental risk posture.
\item \textbf{RUCK} thresholds do not authorize unsafe load progression. Meeting a ruck time under unsafe conditions or with unstable foot/skin integrity does not justify escalation.
\item \textbf{SWIM} and underwater thresholds remain governed. Underwater performance is governance-gated and supervised-only; thresholds do not create permission to practice unsupervised.
\item \textbf{STR} thresholds support capability classification under submax proxy posture and do not override tissue tolerance, technique breakdown, or safety flags.
\item \textbf{AER} thresholds are non-impact conditioning constructs and \textbf{MUST} NOT be treated as interchangeable evidence for \textbf{RUN} tolerance or \textbf{RUN} performance.
\end{itemize}

Decision posture (tier classification logic):

\begin{itemize}
\item Tier classification for progression decisions requires replication: at least 2 of the last 3 \textbf{VALID} tests in the decision window at or above the tier threshold (or at or below, for time metrics).
\item If the last 3 \textbf{VALID} tests straddle tiers, classify to the highest tier that meets replication; otherwise default to the last stable tier that meets replication.
\item If only one \textbf{VALID} test exists in the decision window, the tier classification defaults to \textbf{HOLD} for progression decisions for that metric until replication exists.
\end{itemize}

Versioning note (embedded; governance posture):

\begin{itemize}
\item \textbf{SELECTION-READY} thresholds aligned to Stage 0.3 anchors are frozen as the end-state envelope and \textbf{MUST} NOT move without versioned change control.
\item \textbf{ENTRY} and \textbf{INTERMEDIATE} thresholds anchored to published minimums or competitive equivalents are reference anchors and \textbf{MUST} be verified against current official standards prior to operational use.
\item \textbf{STR} and \textbf{AER} thresholds are governance thresholds not anchored in Stage 0.3 and are versionable via logged change with downstream revalidation once constructs, equipment access, and protocol details are stabilized.
\end{itemize}

\clearpage
\subsection*{2.D.5 Conflict Resolution Rule}
\phantomsection
\addcontentsline{toc}{subsection}{2.D.5 Conflict Resolution Rule}

\begin{itemize}
\item If any downstream phase or gate standard does not map to these tier thresholds by Metric \textbf{ID}, or changes thresholds without versioned change control, that downstream artifact is Draft — Non-Deployable until corrected and re-versioned.
\item Threshold use requires \textbf{VALID} tests. Invalid evidence cannot be used to claim tier attainment.
\end{itemize}

\SubTitleLine

\subsection*{2.D.6 Master Threshold Table Field Definitions and Completion Rules}
\phantomsection
\addcontentsline{toc}{subsection}{2.D.6 Master Threshold Table Field Definitions and Completion Rules}

This section defines completion standards for the master threshold table to ensure every row is auditable and non-ambiguous.

Metric \textbf{ID}

\begin{itemize}
\item Required format: \textbf{MET-2B}-\#\#\# or explicitly noted placeholder \textbf{MET-2B-TBD}-\#\#\#.
\item Completion rule: must match Stage 2.B inventory exactly (placeholders are allowed only under the Assumption Ledger constraints and must be reconciled before Stage 2 is deployable).
\end{itemize}

Metric Name

\begin{itemize}
\item Must match Stage 2.B metric name exactly (or include a crosswalk note if Stage 2.B name differs).
\item Must remain stable for crosswalk joins and phase-row imports.
\end{itemize}

Domain Tag

\begin{itemize}
\item Must match Stage 2.A canonical domain tags: \textbf{RUN}, \textbf{RUCK}, \textbf{SWIM}, \textbf{CAL}, \textbf{STR}, \textbf{BODYCOMP}, \textbf{DURABILITY}, \textbf{RECOVERY}, \textbf{AER}.
\item Must not use synonyms or non-canonical labels.
\end{itemize}

Protocol Card

\begin{itemize}
\item Must be one of:
  \begin{itemize}
  \item \textbf{C1}–\textbf{C18} where protocol exists, or
  \item \textbf{TBD} (2.C) where protocol is not yet drafted, or
  \item N/A only when the metric is a governance-computed compliance statistic (e.g., percent rate) rather than a discrete test procedure.
  \end{itemize}
\item Gate-admissible threshold use requires a drafted protocol for any discrete test metric; where protocol is \textbf{TBD} (2.C), thresholds are placeholders and must be treated as non-gate-admissible until protocol exists.
\end{itemize}

Directionality

\begin{itemize}
\item Must state how the threshold is evaluated:
  \begin{itemize}
  \item Lower is better (time),
  \item Higher is better (reps/seconds),
  \item Pass band (range),
  \item Binary (Y/N).
  \end{itemize}
\item For time thresholds, the mm:ss representation must be interpreted as “at or below” the stated time.
\end{itemize}

Entry Threshold / Intermediate Threshold / Selection-Ready Threshold

\begin{itemize}
\item Must be expressed in the correct unit and directionality for the metric.
\item Must not be based on \textbf{INVALID} evidence.
\item Must not be used without replication when tier classification is used for gate decisions.
\end{itemize}

Notes

\begin{itemize}
\item Required when:
  \begin{itemize}
  \item threshold is derived rather than anchored,
  \item protocol is \textbf{TBD},
  \item comparability breaks are expected,
  \item governance gating controls apply (underwater, disqualifier flags, compliance statistics).
  \end{itemize}
\item Must include derivation basis statement for any derived \textbf{ENTRY}/\textbf{INTERMEDIATE} thresholds and any \textbf{STR}/\textbf{AER} thresholds.
\end{itemize}

Overbuild Preferred

\begin{itemize}
\item Optional performance preference posture only.
\item Must not be interpreted as a gate requirement.
\item Must not be used to justify unsafe escalation or to override validity, durability, or safety posture.
\end{itemize}
\clearpage
\subsection*{2.D.7 Deterministic Tier Classification Algorithm for Gate Use}
\phantomsection
\addcontentsline{toc}{subsection}{2.D.7 Deterministic Tier Classification Algorithm for Gate Use}

Tier classification for progression decisions must be determined by an explicit, repeatable algorithm.

\subsubsection*{2.D.7.1 Pre-checks}
\phantomsection

For a given Metric \textbf{ID}:

1. Evidence availability

\begin{itemize}
\item If there are 0 \textbf{VALID} test events in the decision window → \textbf{UNCLASSIFIED}.
\item If there is 1 \textbf{VALID} test event in the decision window → \textbf{HOLD} (Replication).
\item If there are 2+ \textbf{VALID} test events in the decision window → proceed to replication evaluation.
\end{itemize}

2. Comparability pre-check

\begin{itemize}
\item Confirm the events used are comparable under the protocol’s validity requirements:
  \begin{itemize}
  \item same unit (yards vs meters),
  \item consistent measurement integrity posture (route verified vs treadmill distance basis),
  \item consistent treadmill grade posture where applicable,
  \item load verification posture consistent for rucks,
  \item method consistency posture for trend windows where applicable.
  \end{itemize}
\end{itemize}

If comparability is broken and not stabilized by repeated comparable events → \textbf{HOLD} (Comparability).

3. Disqualifier pre-check

Before tiering performance results, confirm disqualifier controls do not constrain progression use:

\begin{itemize}
\item Any stop-rule event relevant to the metric that compromises safety posture → tier classification may be recorded informationally, but gate use defaults conservative (\textbf{HOLD} or \textbf{MEDICAL} \textbf{HOLD}).
\item Any gait-altering pain flag = Yes within the decision window -> tier classification defaults to \textbf{HOLD} for progression decisions for metrics affected by gait integrity (\textbf{RUCK} and \textbf{RUN} constructs affected).
\item Underwater governance:
  \begin{itemize}
  \item If \textbf{MET-2B-025} is not Y → underwater performance tiering is NOT \textbf{APPLICABLE}; \textbf{MET-2B-026} must not be tiered for gate use.
  \end{itemize}
\end{itemize}

\subsubsection*{2.D.7.2 Replication Evaluation }
\phantomsection

For each tier (\textbf{ENTRY} → \textbf{INTERMEDIATE} → \textbf{SELECTION-READY}):

\begin{itemize}
\item Identify the last 3 \textbf{VALID} test events within the decision window (or fewer if less exist).
\item Count how many of those events meet the tier threshold under directionality:
  \begin{itemize}
  \item time metrics: time ≤ threshold,
  \item rep/hold metrics: value ≥ threshold,
  \item pass-band: value within band,
  \item binary: matches required state.
  \end{itemize}
\item Tier is satisfied for gate use only if at least 2 events satisfy the tier threshold.
\end{itemize}

Classification output:

\begin{itemize}
\item Assign the highest tier satisfied under replication.
\item If none are satisfied, assign \textbf{UNCLASSIFIED} for gate context and treat progression conservatively.
\end{itemize}

\subsubsection*{2.D.7.3 Straddling Tiers}
\phantomsection

If the last 3 \textbf{VALID} events include mixed tier performance:

\begin{itemize}
\item Classify to the highest tier that meets replication.
\item If no tier meets replication but the best single event meets a higher tier:
  \begin{itemize}
  \item record the result,
  \item tier posture remains \textbf{HOLD} (Replication),
  \item do not use the higher tier for progression decisions.
  \end{itemize}
\end{itemize}

\subsubsection*{2.D.7.4 Gate Context Usage}
\phantomsection

\begin{itemize}
\item Tiering can support readiness classification and gate interpretation.
\item Tiering cannot override:
  \begin{itemize}
  \item safety stop rules,
  \item durability disqualifiers,
  \item missing validity condition fields,
  \item missing finalized protocol coverage (placeholder or unresolved protocols cannot be gate-admissible),
  \item Stage 1.F gate evidence rules.
  \end{itemize}
\end{itemize}

\clearpage
\subsection*{2.D.8 Compliance and Disqualifier Metrics: Computation Rules}
\phantomsection
\addcontentsline{toc}{subsection}{2.D.8 Compliance and Disqualifier Metrics: Computation Rules}

These metrics are not “performance tests” but are tiered for governance decision support.

\subsubsection*{2.D.8.1 \textbf{MET-2B-046} Weekly Admissible Session Compliance Percent}
\phantomsection

Definition:

\begin{itemize}
\item Percent of admissible sessions in a week.
\item Admissible session definition aligns to Stage 1.F: \textbf{VALID} plus admissible \textbf{MODIFIED} within cap posture; \textbf{INVALID} excluded.
\end{itemize}

Computation rules:

\begin{itemize}
\item Weekly percent = (admissible sessions ÷ 6) × 100.
\item Decision window percent (two-week window) is the average of the two weekly percents, unless Stage 1.F defines an alternate computation for a specific gate.
\end{itemize}

Minimum data rule:

\begin{itemize}
\item If session tags are missing or were not assigned same day, admissible session count cannot be verified; weekly compliance defaults conservative and may be treated as below threshold for gate decisions.
\end{itemize}

\subsubsection*{2.D.8.2 \textbf{MET-2B-045} Cardio Minimum Standard Met}
\phantomsection

Definition:

\begin{itemize}
\item Y/N per session indicating whether the Cardio minimum standard was met and documented.
\end{itemize}

Computation rules:

\begin{itemize}
\item Cardio met requires:
  \begin{itemize}
  \item modality recorded,
  \item minutes recorded,
  \item continuity posture satisfied,
  \item intensity anchor recorded (RPE + talk-test cue),
  \item and the underlying session is not invalidated by missing required elements.
  \end{itemize}
\end{itemize}

If Cardio minimum is missed:

\begin{itemize}
\item The session is invalidated per Stage 1.F posture and cannot be counted admissible for gate computations.
\item Therefore the decision-window cardio-compliance rate derived from \textbf{MET-2B-045} is both:
  \begin{itemize}
  \item a direct compliance metric, and
  \item an indirect driver of admissible session count.
  \end{itemize}
\end{itemize}

\subsubsection*{2.D.8.3 \textbf{MET-2B-TBD-001} Weekly \textbf{MODIFIED} Session Count (Draft — Non-Deployable until reconciled into Stage 2.B)}
\phantomsection

Definition:

\begin{itemize}
\item Count of \textbf{MODIFIED} sessions per week where the modification is not driven by clinician-imposed restriction posture or stop-rule events requiring conservative substitution.
\end{itemize}

Interpretation posture:

\begin{itemize}
\item This metric expresses drift and substitution frequency.
\item If this count exceeds thresholds repeatedly, gate posture defaults to conservative even if performance thresholds are improving.
\end{itemize}

\subsubsection*{2.D.8.4 \textbf{MET-2B-TBD-002} Weekly \textbf{INVALID} Session Count (Draft — Non-Deployable until reconciled into Stage 2.B)}
\phantomsection

Definition:

\begin{itemize}
\item Count of \textbf{INVALID} sessions per week.
\end{itemize}

Interpretation posture:

\begin{itemize}
\item \textbf{INVALID} sessions are non-admissible.
\item Any sustained pattern of \textbf{INVALID} sessions forces conservative gate posture:
  \begin{itemize}
  \item \textbf{HOLD} by default,
  \item \textbf{REGRESS} if tolerance failure or repeated invalidity persists,
  \item \textbf{MEDICAL} \textbf{HOLD} if invalidity is driven by stop-rule categories.
  \end{itemize}
\end{itemize}

\subsubsection*{2.D.8.5 \textbf{MET-2B-036} Gait-Altering Pain Flag}
\phantomsection 

Interpretation posture:

\begin{itemize}
\item For progression decisions:
  \begin{itemize}
  \item Any Yes within the decision window forces \textbf{HOLD} for tier-based progression decisions in impact/load domains.
  \end{itemize}
\item For test events:
  \begin{itemize}
  \item The flag \textbf{MUST} be No during any \textbf{VALID} test event for tier-classification use at \textbf{INTERMEDIATE} and \textbf{SELECTION-READY}.
  \end{itemize}
\end{itemize}

\clearpage
\subsection*{2.D.9 Protocol Coverage Gaps and Gate-Admissibility Handling}
\phantomsection
\addcontentsline{toc}{subsection}{2.D.9 Protocol Coverage Gaps and Gate-Admissibility Handling}

Protocol coverage gaps must be treated conservatively.

\begin{itemize}
\item If Protocol Card = \textbf{TBD} (2.C):
  \begin{itemize}
  \item Thresholds are placeholders only.
  \item The metric cannot be treated as gate-admissible until the protocol exists and defines:
    \begin{itemize}
    \item validity fields,
    \item measurement integrity rules,
    \item stop-rule overlays,
    \item and invalidation triggers.
    \end{itemize}
  \end{itemize}
\end{itemize}

This applies explicitly and only to metrics whose protocol card remains unresolved or placeholder-incomplete:

\begin{itemize}
\item \textbf{MET-2B-032} through \textbf{MET-2B-034} (\textbf{AER} time trials) remain supporting only until \textbf{C18} is fully finalized and non-placeholder.
\end{itemize}
\clearpage
\subsection*{2.D.10 Change Control and Revalidation Posture for Thresholds}
\phantomsection
\addcontentsline{toc}{subsection}{2.D.10 Change Control and Revalidation Posture for Thresholds}

Thresholds are governance-critical. Changes must be controlled.

\begin{itemize}
\item Any change to \textbf{SELECTION-READY} thresholds aligned to Stage 0.3 anchors is treated as a constraint-level change and requires explicit change control and downstream revalidation.
\item Any change to \textbf{ENTRY}/\textbf{INTERMEDIATE} thresholds must include:
  \begin{itemize}
  \item rationale,
  \item derivation basis,
  \item comparability impact statement (if any),
  \item and downstream revalidation of affected phase rows and gate logic.
  \end{itemize}
\item Any change to a metric’s protocol reference (C\#) is a comparability-impacting change and requires revalidation of:
  \begin{itemize}
  \item the protocol card,
  \item affected thresholds,
  \item and any downstream consumers referencing the metric.
  \end{itemize}
\end{itemize}

\subsection*{2.D.11 Acceptance Criteria for Stage 2.D}
\phantomsection
\addcontentsline{toc}{subsection}{2.D.11 Acceptance Criteria for Stage 2.D}

Stage 2.D is deployable only when:

\begin{itemize}
\item Every Gate-Critical metric in Stage 2.B has:
  \begin{itemize}
  \item tier thresholds for \textbf{ENTRY}, \textbf{INTERMEDIATE}, \textbf{SELECTION-READY} (or explicit N/A posture where governance locks the domain, e.g., underwater).
  \end{itemize}
\item Directionality is explicit and consistent with units and protocol constructs.
\item Replication rule is explicitly applied and does not conflict with Stage 1.F gate evidence posture.
\item Any placeholder metrics (\textbf{MET-2B-TBD}-\#\#\#) are either:
  \begin{itemize}
  \item reconciled into Stage 2.B with permanent IDs, or
  \item explicitly marked as Draft — Non-Deployable dependencies until reconciled.
  \end{itemize}
\item Any metrics with Protocol Card = \textbf{TBD} (2.C) are treated as non-gate-admissible until protocol coverage is drafted and validated.
\item No thresholds imply permission to violate safety posture, durability disqualifiers, or water governance boundaries.
\end{itemize}

\end{document}